\chapter{Output reference}\label{ch:outputs}

\section{Phase 1}
\begin{description}
\item[\file{resonanceslist}] One line per resonant term: type (B3, D4, E4,
  G2, \ldots), modes, coupling, detuning and their ratio.
\item[\file{dunhamparams}]\label{sec:dunham} The deperturbed expansion
  $E(n)=E_0+\sum_i w_i(n_i+\frac12)+\sum_{i\le j}x_{ij}(n_i+\frac12)(n_j+\frac12)$,
  with each pair counted once, as in Gaussian's X matrix. It gives $E_0$,
  the zero-point energy, $\omega_i$, $w_i$, the deperturbed fundamentals
  $\nu_i$, and all the $x_{ij}$. In the theory notes the sum runs over
  ordered pairs, so $x_{ij}=2\chi_{ij}$ for $i\ne j$ \notes{§6.6}.
\item[\file{bandlist}] \devonly{} All bands with their dipole derivative, harmonic
  wavenumber, intensity estimate and direction (\texttt{both}, \texttt{up},
  \texttt{down}, \texttt{off}).
\item[\file{resonancecandidates}] \devonly{} All candidate resonances with their
  metrics, decision and reason, grouped by change of quantum numbers. Terms
  that are resonant only because another term with the same change is
  (they are ``carried along'') are marked \texttt{carried}, so that the
  terms with decision~1 are exactly those of \file{resonanceslist}. The
  last column is the override in force (\texttt{on}, \texttt{off}) or
  \texttt{auto}.
\item[\file{*.shifted}, \file{gradient.shifted}, \file{shiftgeometry.out}]
  With \opt{shiftgeometry}.
\item[\file{dumpedparams}] Binary, read by phases 2a and 2b. It must come
  from the same version of phase 1.
\end{description}
\section{Phase 2a}
\begin{description}
\item[\file{dir\_nNNNNN}] One directory per starting state, with the
  polyad structures and Hamiltonians.
\item[\file{processedquantas}, \file{processeddirs}] The starting states
  and their directories.
\item[\file{polyadsizes}] The polyad sizes.
\end{description}
\section{Phase 2b}
\begin{description}
\item[\file{hist\_abs\_0}, \file{\_1}, \file{\_2}, \file{hist\_ems\_*}]
  Absorption and emission histograms, one file per Cartesian component of
  the dipole, binned in wavenumber and in the harmonic energy of the
  starting state. The total spectrum is the sum of the three files.
\item[\file{hist\_anh\_*}] The same, binned in the anharmonic energy of the
  starting state.
\item[\file{hist\_energy}, \file{hist\_anhenergy}] Number of starting
  states per energy bin.
\item[\file{proclog\_n*}] Logs of each process, with the truncation
  estimates when \opt{residualcheck} is on.
\end{description}
% TODO: units and normalisation of the histograms; script to convert them
% into a spectrum in km/mol.
